\documentclass[UTF8,a4paper,11pt,openany]{ctexbook}
\usepackage{../common/statlearnbook}
\SLSetBookMetadata{第 5 册}{统计学习方法讲义 v2.6.0}{采样方法、主题模型与图排序}
\begin{document}
\setcounter{page}{769}
\setcounter{chapter}{32}
\setcounter{figure}{6}
\pagestyle{headings}
\chaptermark{马尔可夫链蒙特卡罗法}
\sectionmark{收敛与样本诊断}
每个坐标更新核 $K_j$ 都保持目标分布 $\pi$，但核的复合与混合具有不同的可逆性。图\ref{fig:V5-C03-componentwise-sweep}区分“一个坐标更新”与“完成一轮扫描”，并标出系统扫描和随机扫描在可逆性上的差异。
% v2.3.1 figure UID: FIG-P605-01
% v2.6.0 reverse redraw for FIG-P605-01: aligned semantic rows and collision-free spacing.
\tikzset{slfig-FIG-P605-01/.style={font=\fontsize{9.2pt}{11.0pt}\selectfont,every path/.append style={line cap=round,line join=round}}}
\begin{figure}[htbp]
\centering
\begin{tikzpicture}[slfig-FIG-P605-01,>=Stealth,every node/.style={font=\fontsize{9.2pt}{11.0pt}\selectfont,align=center},
  kernel/.style={draw=SLBlue,rounded corners=2pt,fill=white,minimum width=14mm,minimum height=9mm,line width=.72pt},
  choice/.style={diamond,aspect=2.1,draw=SLGold,fill=SLGold!8,text width=18mm,inner sep=1pt,line width=.78pt},
  flow/.style={-{Stealth[length=2.0mm]},draw=SLBlue,line width=.95pt},
  branch/.style={-{Stealth[length=1.8mm]},draw=SLGray,line width=.70pt}]
  \draw[draw=SLRuleGray,line width=.55pt,rounded corners=3pt,fill=SLSoftGray] (-7.4,-2.80) rectangle (-.55,2.55);
  \draw[draw=SLRuleGray,line width=.55pt,rounded corners=3pt,fill=SLSoftGray] (.55,-2.80) rectangle (7.4,2.55);
  \node[font=\fontsize{9.8pt}{11.8pt}\selectfont\bfseries,text=SLBlue] at (-3.98,2.02) {系统扫描：固定复合};
  \node[kernel] (k1) at (-6.20,.78) {$K_1$};
  \node[kernel] (k2) at (-4.40,.78) {$K_2$};
  \node (dots) at (-3.02,.78) {$\cdots$};
  \node[kernel] (kd) at (-1.55,.78) {$K_d$};
  \draw[flow] (k1)--(k2); \draw[flow] (k2)--(dots); \draw[flow] (dots)--(kd);
  \node at (-3.98,-.35) {$K_{\rm sys}=K_1K_2\cdots K_d$};
  \node[draw=SLGray,line width=.58pt,rounded corners=2pt,fill=white,minimum width=56mm,minimum height=7mm] at (-3.98,-2.05)
    {固定次序复合：通常不保证可逆};

  \node[font=\fontsize{9.8pt}{11.8pt}\selectfont\bfseries,text=SLBlue] at (3.98,2.02) {随机扫描：先抽坐标};
  \node[choice] (pick) at (3.98,1.05) {$J\sim\omega$};
  \node[kernel] (r1) at (1.55,-.30) {$K_1$};
  \node[kernel] (rj) at (3.98,-.30) {$K_j$};
  \node[kernel] (rd) at (6.35,-.30) {$K_d$};
  \draw[branch] (pick)--(r1); \draw[branch] (pick)--(rj); \draw[branch] (pick)--(rd);
  \node at (3.98,-1.20) {$K_{\rm rand}=\sum_{j=1}^{d}\omega_jK_j$};
  \node[draw=SLBlue,line width=.72pt,rounded corners=2pt,fill=white,
    text width=52mm,minimum height=10mm,inner xsep=3mm] at (3.98,-2.20)
    {若各 $K_j$ 关于 $\pi$ 可逆\\固定权重混合保持可逆};
\end{tikzpicture}
\caption{分量MH的系统扫描与随机扫描。固定顺序的核复合通常不保证可逆；可逆坐标核的固定权重随机混合仍保持可逆。}\label{fig:V5-C03-componentwise-sweep}
\end{figure}

\begin{theorem}[有限状态 MH 的正确性条件]
设状态空间有限且 $\pi(x)>0$。若 MH 核不可约且非周期，则 $\pi$ 是唯一平稳分布，并且链从任意初始分布收敛到 $\pi$。
\end{theorem}
\end{document}
